\documentclass[aspectratio=169]{beamer}

\usetheme{Madrid}
\usecolortheme{default}

\usepackage{graphicx}
\usepackage{amsmath}
\usepackage{amssymb}
\usepackage{dsfont}
\usepackage{booktabs}
\usepackage{xcolor}

%%% Notation (same as report) %%%
\newcommand{\states}{\ensuremath{\mathcal{S}}}
\newcommand{\actions}{\ensuremath{\mathcal{A}}}
\newcommand{\policy}{\ensuremath{\pi}}
\newcommand{\actionval}{\ensuremath{q}}
\newcommand{\optaval}{\ensuremath{q_*}}
\newcommand{\reward}{\ensuremath{\mathcal{R}}}
\newcommand{\transition}{\ensuremath{\mathcal{P}}}
\newcommand{\discount}{\ensuremath{\gamma}}

% -----------------------------

\title{Warehouse Navigation via Reinforcement Learning}
\subtitle{Reinforcement Learning and Decision Making Under Uncertainty}
\author{NAME 1 \and NAME 2 \and NAME 3}
\institute{University Name}
\date{Spring 2026}

% -----------------------------
% ============================================================
% TIMING GUIDE  (7 min talk + 3 min Q&A)
%   Slide 1  — Title              ~0:00–0:20  (20 s)
%   Slide 2  — Problem & MDP      ~0:20–1:30  (70 s)
%   Slide 3  — Environment        ~1:30–2:30  (60 s)
%   Slide 4  — Algorithms         ~2:30–4:00  (90 s)
%   Slide 5  — Results            ~4:00–5:30  (90 s)
%   Slide 6  — Robustness         ~5:30–6:30  (60 s)
%   Slide 7  — Conclusion         ~6:30–7:00  (30 s)
% ============================================================

\begin{document}

% --- Slide 1 (~20 s) ---
\begin{frame}
  \titlepage
\end{frame}

% --- Slide 2 (~70 s) ---
\begin{frame}{Problem, MDP \& Approach}
  % What is the problem, how is it modelled as an MDP, which algorithms were used?
  \begin{columns}
    \column{0.5\textwidth}
      \textbf{Problem}
      \begin{itemize}
        \item Fragestellung/Problemstellung
        \item
      \end{itemize}
      \textbf{MDP: } $(\states, \actions, \transition, \reward, \discount)$
      \begin{itemize}
        \item $\states$ —
        \item $\actions$ —
        \item $\reward$ —
        \item $\discount$ —
      \end{itemize}
    \column{0.5\textwidth}
      \textbf{Algorithms}
      \begin{itemize}
        \item Q-Learning
        \item SARSA
        \item DQN
        \item Value Iteration
      \end{itemize}
  \end{columns}
\end{frame}

% --- Slide 3 (~60 s) ---
\begin{frame}{Environment}
  % Grid layout figure on the right; key properties on the left.
  \begin{columns}
    \column{0.5\textwidth}
      \begin{itemize}
        \item Grid size:
        \item Obstacles:
        \item Start / Goal:
        \item Stochasticity:
        \item Episode ends when:
      \end{itemize}
    \column{0.5\textwidth}
      % \includegraphics[width=\textwidth]{figures/warehouse_env.png}
      [FIGURE]
  \end{columns}
\end{frame}

% --- Slide 4 (~90 s) ---
\begin{frame}{Algorithms \& Key Equations}
  % Show the update rule for each algorithm side by side or stacked.
  \textbf{Q-Learning / SARSA update:}
  \[
    \actionval(s_t,a_t) \leftarrow \actionval(s_t,a_t)
    + \alpha\bigl[r_t + \discount\, \actionval(s_{t+1},a') - \actionval(s_t,a_t)\bigr]
  \]
  % Q-Learning: a' = argmax; SARSA: a' ~ policy

  \textbf{DQN loss:}
  \[
    \mathcal{L}(\theta) = \mathbb{E}\!\left[
      \bigl(r + \discount\max_{a'}\actionval(s',a';\theta^{-})
      - \actionval(s,a;\theta)\bigr)^2\right]
  \]

  \begin{table}
    \centering
    \small
    \begin{tabular}{lccc}
      \toprule
      Parameter & Q-Learning & SARSA & DQN \\
      \midrule
      $\alpha$          & & & \\
      $\discount$       & & & \\
      $\varepsilon$     & & & \\
      Episodes          & & & \\
      \bottomrule
    \end{tabular}
  \end{table}
\end{frame}

% --- Slide 5 (~60 s) ---
\begin{frame}{Deep Q-Learning}
  \begin{columns}

    \column{0.52\textwidth}

      \begin{itemize}
        \item DQN applied to warehouse MDP
        \item Initial training produced only negative rewards
      \end{itemize}

      \vspace{0.3cm}

      \textbf{Tried Improvements}
      \begin{itemize}
        \item Changed state encoding
          \begin{itemize}
            \item Grid-based state representation
            \item Channels: obstacles, packages, storage, robot, carrying
          \end{itemize}

        \item Implemented action masking

        \item Tuned network hyperparameters
          \begin{itemize}
            \item Slow $\varepsilon$-decay
            \item Learning-rate scheduling
          \end{itemize}
      \end{itemize}

    \column{0.48\textwidth}

      \textbf{Main Challenges}
      \begin{itemize}
        \item Agent did not reliably learn
              when pickup / delivery actions are beneficial

        \item Some runs achieved deliveries

        \item Other runs converged to
              negative rewards only
      \end{itemize}

  \end{columns}
\end{frame}

% --- Slide 6 (~30 s) ---
\begin{frame}{Deep Q-Learning Results}
  \begin{columns}
    \column{0.5\textwidth}
      \begin{center}
        \includegraphics[width=\textwidth]{report/figures/warehouse_rewards_delivered.png}
      \end{center}
    \column{0.5\textwidth}
      \begin{center}
        \includegraphics[width=\textwidth]{report/figures/warehouse_rewards_delivered_2.png}
      \end{center}
  \end{columns}
\end{frame}

% --- Slide 7 (~60 s) ---
\begin{frame}{Robustness}
  % What perturbations were tested? What held up and what broke?
  \begin{itemize}
    \item Stochasticity ($p$ varied):
    \item Unseen layouts:
    \item Reward perturbations:
    \item Hyperparameter sensitivity:
  \end{itemize}
\end{frame}

% --- Slide 7 (~30 s) ---
\begin{frame}{Conclusion}
  % One bullet per point — keep it tight.
  \textbf{Summary}
  \begin{itemize}
    \item
  \end{itemize}
  \textbf{Limitations \& Future Work}
  \begin{itemize}
    \item
    \item
  \end{itemize}
\end{frame}

\end{document}
